% =====================================================================
%  tr04a.tex —— 原文 pp.53--62（译者：tr-04）
%  内容：§7 第六讲（Adams 运算的应用：Hopf 不变量为 1 的映射的不存在性）；
%        §8 解析 K 理论（K-同调）总起；§9 第一讲（泛函分析的一些预备，
%        Fredholm 算子的指标）
%  原文无编号公式：展示式一律 \[ … \]；定理类号数照原文手工给出。
% =====================================================================
% ===== 原文 pp.53--62 =====
% [原文 p.53]
\section{第六讲（Adams 运算的应用：Hopf 不变量为 1 的映射的不存在性）}

\emph{作为 Adams 运算的一个应用，我们将证明 Hopf 不变量为 1 的映射不存在。}

令
\[ S^{4n-1}\xra{f}S^{2n}\to S^{2n}\cup_f C^{4n}=X. \]

我们算出
\[ \wt K^0(X)=\Z\oplus\Z, \]
其中用到对 $(X,S^{2n})$ 的正合列以及 $X/S^{2n}=S^{4n}$。

取若干生成元 $p:X\to X/S^{2n}=S^{4n}$ 与 $p^*:\wt K^0(S^{4n})\to\wt K^0(X)$，其中 $u_{4n}$ 是 $\wt K^0(S^{4n})=\Z$ 的生成元，而 $y=p^*(u_{4n})\in\wt K^0(X)$。

接下来我们选取元素 $x\in\wt K^0(X)$，使得
\[ i^*(x)=u_{2n}, \]
其中
\[ i^*:\wt K^0(X)\to\wt K^0(S^{2n}), \]
并且 $u_{2n}\in\wt K^0(S^{2n})$；我们考虑
\[ x^2:\quad i^*(x^2)=0. \]

\para{定义.} 映射 $f$ 的\emph{Hopf 不变量}是整数 $\lambda_f$，它由方程
\[ x^2=\lambda_f y \]
定义。

可以证明，$\lambda_f$ 与 $x$ 的选取无关。

% [Tutor Note: p.53 这里的 $C^{4n}$ 不是复射影空间，而是 $4n$ 维胞腔（即 $S^{4n-1}$ 上的锥）；$X=S^{2n}\cup_f C^{4n}$ 是沿 $f$ 把 $4n$ 胞腔粘到 $S^{2n}$ 上所得的胞腔复形。核对 $X/S^{2n}=S^{4n}$：把 $S^{2n}$ 压成一点后只剩一个 $4n$ 球面。$\wt K^0(X)=\Z\oplus\Z$ 来自对 $(X,S^{2n})$ 的正合列，它分裂为 $0\to\wt K^0(S^{4n})\to\wt K^0(X)\to\wt K^0(S^{2n})\to 0$，故 $\Ker i^*=\Z y$。原文说 "Pick some generators $p$..." 措辞不严谨：$p$ 是商映射而不是生成元，真正的第二个生成元是 $p^*(u_{4n})=y$。]

% [Tutor Note: p.53 $\lambda_f$ 之所以能定义，是因为 $i^*(x^2)=i^*(x)^2=u_{2n}^2=0$（$K^0(S^{2n})\cong\Z[t]/(t^2)$，约化部分平方为零），于是 $x^2\in\Ker i^*=\Z y$，可写成 $\lambda_f y$。这正是「Hopf 不变量」在 K 理论中的翻译：它衡量 $x^2$ 与 $y$ 的倍数关系；原文的 $x^2:i^*(x^2)=0$ 就是为此做的准备。]

% [Tutor Note: p.53 「可以证明 $\lambda_f$ 与 $x$ 的选取无关」的补证：$x$ 的任一其他选取形如 $x'=x+cy$（$c\in\Z$），因为 $i^*(x'-x)=0$ 而 $\Ker i^*=\Z y$。于是 $x'^2=x^2+2cxy+c^2y^2$；在 K 理论的 Atiyah--Hirzebruch 过滤中 $x\in F^{2n}$、$y\in F^{4n}$（$y$ 在 $2n$ 骨架上为零），故 $xy\in F^{6n}=0$、$y^2=0$，从而 $x'^2=x^2$。原文省略了这段论证。]

于是有如下结果：

\begin{thm}{1} 若 $\lambda_f$ 为奇数，则 $n=1,2,4$。 \end{thm}

% [原文 p.54]
\begin{proof}
回忆
\[ \wt K^0(X)=\Z\oplus\Z, \]
其生成元记作 $x$ 与 $y$（每个 $\Z$ 对应一个）。于是：
\[ \Psi^2(x)=2^n x+ay \]
以及
\[ \Psi^3(x)=3^n x+by, \]
而
\[ \Psi^k(y)=\Psi^k(p^*(u_{4n}))=p^*(\Psi^k(u_{4n}))=p^*(k^{2n}u_{4n})=k^{2n}y, \]
其中
\[ \Psi^2(x)\mapsto\Psi^2(u_{2n})=2^n u_{2n} \]
以及
\[ \Psi^3(x)\mapsto\Psi^3(u_{2n})=3^n u_{2n}. \]
此外我们计算：
\[ \Psi^6(x)=\Psi^3(\Psi^2(x))=6^n x+(2^n b+3^{2n}a)y \]
\[ \Psi^6(x)=\Psi^2(\Psi^3(x))=6^n x+(2^{2n}b+3^n a)y, \]
故
\[ 2^n b+3^{2n}a=2^{2n}b+3^n a\iff 2^n(2^n-1)b=3^n(3^n-1)a. \]
但
\[ \Psi^2(x)=x^2\bmod 2=\lambda_f y\bmod 2=y\bmod 2\ \Rightarrow\ a\ \text{为奇数}. \]
由此，从
\[ 2^n(2^n-1)b=3^n(3^n-1)a \]
推出 $2^n$ 整除 $3^n-1$。直接计算可得 $n=1,2,4$。
\end{proof}

有三类经典的 Hopf 映射：

% [Tutor Note: p.54 两处关键事实：① $\Psi^k(u_{4n})=k^{2n}u_{4n}$，因为 $u_{4n}$ 是 $\wt K^0(S^{4n})=\Z$ 的生成元，而 $\Psi^k$ 在 $\wt K^0(S^{2m})$ 上就是乘 $k^m$（$S^{4n}$ 的生成元是 Bott 类的 $2n$ 次幂）；② 链式等式的第一步 $\Psi^2(x)=x^2\bmod 2$ 来自 λ 环恒等式 $\Psi^2(x)=x^2-2\lambda^2(x)$。原文的 $\Psi^2(x)\mapsto\Psi^2(u_{2n})$ 是 $i^*\Psi^2(x)=\Psi^2(i^*x)$ 的压缩写法（省掉了 $i^*$）。另，同页 "taht" 是 "that" 的排印颠倒，正文照译。]

% [原文 p.55]
\[ S^3\to S^2,\quad S(\C^2)=S^3\to S^2=\C P^1,\quad (z_1,z_2)\mapsto[z_1,z_2] \]
\[ S^7\to S^4,\quad S(\HH^2)=S^7\to S^4=\HH P^1,\quad (q_1,q_2)\mapsto[q_1,q_2] \]
\[ S^{15}\to S^8,\quad S(\mathbb{O}^2)=S^{15}\to S^8=\mathbb{O}P^1,\quad (k_1,k_2)\mapsto[k_1,k_2] \]

其中 $\HH$ 表示四元数，而 $\mathbb{O}$ 表示 Caley 数（即八元数）。

\emph{由上面的定理 1 可知，不存在其他这样的映射。}

一个类似的问题是：\emph{在数列}
\[ \text{实数}\to\text{复数}\to\text{四元数}\to\text{Caley 数} \]
\emph{中，还有别的数吗？}

\emph{答案也是否定的。}

我们以如下注记结束本讲：圆周 $S^1$（在拓扑上）就是群 $(U(1)$。

球面 $S^3$（在拓扑上）是单位四元数构成的群。

球面 $S^7$ 在拓扑上是单位 Caley 数构成的群。

空间 $1\times S^7\subset S^7\times S^7\to S^7$ 也是一个「群」，即它是一个 H-空间；而 $S^7\times 1\subset S^7\times S^7\to S^7$ 亦然。换言之，我们想知道球面 $S^i$ 是否为\emph{可平行化空间}，即它们的切丛是否平凡。

% [Tutor Note: p.55 原文 "Caley numbers" 是 Cayley 数（八元数）的拼写讹误（标准作 Cayley），正文照原文音译并在括号内保留「即八元数」；$\mathbb{O}$ 在本文中表示八元数，勿与正交群 $\OO$ 混淆。同页 "there no other such maps" 缺系动词 are，"the group (U(1)" 多了一个左括号（应为 $U(1)$）——均照原样翻译。又：$S^1,S^3,S^7$ 是仅有的可平行化球面（Bott--Milnor），这正呼应本讲末句「$S^i$ 是否可平行化」；但 $S^7$ 上的八元数乘法不结合，故它只是 H-空间而非拓扑群。]

% [原文 p.56]
\section{解析 K 理论（K-同调）}

解析 K 理论的一个动机可能是这个问题：「\emph{无穷大减无穷大是什么}？」此外，\emph{Fredholm} 算子的\emph{指标}是数学中最有用的定义之一，因为它催生了各种\emph{指标定理}——可以说是 20 世纪后半叶数学中最核心的结果。与此同时，$C^*$-代数的 K 理论是 A. Connes 的\emph{非交换几何}中的一个重要组成部分。

\section{第一讲（泛函分析的一些预备，Fredholm 算子的指标）}

由拓扑 K 理论我们知道 $K_0(*)=\Z$，即 K 理论不满足维数公理，因此它是广义的（或「异域的」）同调理论。

设 $V_0,V_1$ 是两个有限维向量空间（取自某个域 $\F$，通常是复数域）。计算它们的维数之差是直接的，我们把它记作 $[V_0]-[V_1]$。（为简单起见，我们常写 $[V]$ 而不写 $\dim V$。）

设 $T:V_0\to V_1$ 是线性算子。我们展开：
\[ 0\to\Ker T\to V_0\xra{T}V_1\to\coker T\to 0, \]
其中
\[ \coker T=V_1/\im T=V_1/T(V_0)=(\im T)^\perp. \]

\begin{thm}{1（线性代数）}
\[ [\Ker T]-[V_0]+[V_1]-[\coker T]=0. \]
\end{thm}

（该定理对任意有限维向量空间的短正合列都成立。）

% [Tutor Note: p.56 §8 只有这一段总起，没有定义或定理，可看作解析 K 理论部分（§8--§14）的导言：它把动机定在「$\infty-\infty$ 是什么」，并把 $C^*$-代数的 K 理论与 A. Connes 的非交换几何联系起来。$K_0(*)=\Z$ 正是 K 理论不满足维数公理（Eilenberg--Steenrod 公理之一）的体现。]

% [Tutor Note: p.56 「无穷大减无穷大是什么」正是本讲（§9）的出发点：对 Fredholm 算子 $T$，$\dim\Ker T-\dim\coker T$ 是 $\infty-\infty$ 的一个可计算版本，而且（如 p.61 的移位算子所示）它更多取决于空间而不是算子。]

% [Tutor Note: p.56 标题里的 "K-Homology" 是「解析 K 理论」的别名：与拓扑 K 理论（广义上同调）相对，这一支从 $C^*$-代数、Fredholm 算子与扩张出发，得到的是同调式的理论，两者经由 Kasparov 的 KK 理论配对。本书解析部分的 §13、§14 分别讲 $C^*$-代数的 K 理论与 K-同调。]

% [Tutor Note: p.56 $\coker T=V_1/\im T=V_1/T(V_0)=(\im T)^\perp$ 的最后一个等号只在有限维或 $\im T$ 闭时成立（此时 $V_1=\im T\oplus(\im T)^\perp$）。另外注意本讲（§9）的定理编号：原文在 p.56 与 p.59 各有一个「定理 1」，同一节内重号，译文照原样保留，不作重排。]

% [原文 p.57]
我们可以把定理中的关系改写成更方便的形式：
\[ [V_0]-[V_1]=[\Ker T]-[\coker T]. \]

如果我们试图把上述关系推广到\emph{无穷维}向量空间（例如 Hilbert 空间），就会看到左端没有意义，因为它给出 $\infty-\infty$；但右端在某些有利情形下可能给出有意义的结果（只要算子是一个「好」算子）。因此，差 $\infty-\infty$ 在某些情形下可以给出有限的结果。

在本章中，除非另有明确说明，所有算子都假定是线性的。

我们简要回忆 Hilbert 空间的定义：

\begin{defn}{1} Hilbert 空间 $H$ 是一个带内积的无限维完备（复）向量空间。完备是指：$H$ 中每个满足 $\sum\|x_n\|<\infty$ 的序列 $(x_n)$ 都在 $H$ 中收敛。 \end{defn}

\begin{defn}{2} 线性映射（线性算子）$T:H_0\to H_1$ 称为\emph{有界}的，如果存在某个正实数 $c$，使得
\[ \|Tx\|\le c\|x\|,\quad \forall x\in H. \]
\end{defn}

这种正实数中「最好」的那个称为 $T$ 的\emph{范数}，记作 $\|T\|$。特别地，可以定义
\[ \|T\|=\sup\langle Tx,y\rangle, \]
其中 $\|x\|=\|y\|=1$。

所有有界线性算子 $T:H\to H$ 的集合记作 $B(H)$。

给定 Hilbert 空间，我们可以定义\emph{Hilbert 基}的概念；关于上述 Hilbert 空间 $H_0,H_1$ 的某些 Hilbert 基，每个（有界）线性算子都可以用一个 $\infty\times\infty$ 矩阵表示（基本上采用与用矩阵表示有限维向量空间之间的线性映射相同的做法）。反过来，给定一个 $\infty\times\infty$ 矩阵，

% [Tutor Note: p.57 定义 1 用「$\sum\|x_n\|<\infty$ 的级数收敛」刻画完备性；在赋范空间中它等价于通常的 Cauchy 完备性，但原文未说明。定义 2 的范数公式照原文写为 $\|T\|=\sup\langle Tx,y\rangle$，严格应写 $\|T\|=\sup_{\|x\|=\|y\|=1}|\langle Tx,y\rangle|$（取绝对值）。同页 "linear opearator" 为 "linear operator" 的排印讹误，正文照译。]

% [原文 p.58]
要判断它是否定义了一个有界线性算子就困难得多。

然而我们知道：一个 $\infty\times\infty$ 矩阵确实是某个有界线性算子的表示，如果（但不仅如果）它只有有限多个非零元素。在这种情形下它称为\emph{有限秩算子}。换言之，有限秩算子是有界的（但反过来不一定成立）。

\begin{defn}{3} 线性算子 $T$ 称为\emph{紧}的，如果 $\forall\epsilon>0$，存在有限秩算子 $F$ 使得
\[ \|T-F\|<\epsilon. \]
\end{defn}

所有紧算子 $T:H\to H$ 的集合记作 $K(H)$。

$B(H)$ 具有自然的 $C^*$-代数结构，而 $K(H)$ 是 $B(H)$ 中的一个理想（闭的双边 $*$-理想），因此可以作商
\[ Q(H)=B(H)/K(H), \]
它是另一个 $C^*$-代数，称为 \emph{Calkin 代数}（另见后文第五讲）。

例如，恒等算子 $I:H\to H$ \textbf{不是}紧的。

\begin{prop}{1} 设 $B(H)$ 表示 $H$ 到自身的所有有界线性算子的集合。$B(H)$ 中所有可逆算子构成的子集是开集。 \end{prop}

\begin{proof}
只需证明恒等算子 $I$ 是可逆算子集合的内点，因为若 $\|A\|<1$，则 $(I-A)^{-1}=I+A+A^2+\cdots$。
\end{proof}

\begin{defn}{4} 两个 Hilbert 空间之间的有界线性算子 $T:H_0\to H_1$ 称为 \emph{Fredholm} 的，如果它的核 $\Ker T$ 是有限维的，它的像 $\im T$ 是 $H_1$ 的闭子空间，并且商 $H_1/\im T$ 也是有限维的。 \end{defn}

% [Tutor Note: p.58 命题 1 的证明只处理了 $A$ 接近 $I$ 的情形；补出一般情形：若 $T$ 可逆且 $\|T-A\|<1/\|T^{-1}\|$，则 $\|T^{-1}(T-A)\|<1$，由 Neumann 级数 $A=T(I-T^{-1}(T-A))$ 可逆，故可逆算子集合是开集。另 "internal point" 是 "interior point"（内点）的旧式说法。]

% [原文 p.59]
有下述基本结果：

\begin{thm}{1（Atkinson）} 下列各项等价：
\begin{enumerate}[label=\textbf{\alph*.},leftmargin=2.2em,itemsep=3pt]
  \item $T$ 是 Fredholm 的。
  \item $T$ 模紧算子可逆，即存在算子 $S:H_1\to H_0$，使得 $ST-I$ 与 $TS-I$ 分别在 $H_0$ 与 $H_1$ 上是紧的。换言之，「模紧算子可逆」这一说法是指：$T$ 在标准投影 $\pi:B(H)\to Q(H)$ 下的像 $\pi(T)$ 是可逆的。
\end{enumerate}
\end{thm}

\begin{proof}
\textbf{a.} $\Rightarrow$ \textbf{b.}
$H_0$ 是子空间 $(\Ker T)^\perp$ 与 $\Ker T$ 的直和，而 $H_1$ 是 $\im T$ 与 $(\im T)^\perp\simeq H_1/\im T$ 的直和，因此存在线性映射 $(\Ker T)^\perp\to\im T$。借助泛函分析中的闭图像定理（见附录）推出，上述连续双射有连续逆 $S:(\im T)\to(\im T)^\perp$。于是 $ST-I$ 与 $TS-I$ 都是有限秩算子。

\textbf{b.} $\Rightarrow$ \textbf{a.}
设存在算子 $S$，使得 $ST-I$ 与 $TS-I$ 都是紧的。我们把第一个限制到 $\Ker T$ 上，于是
\[ (ST-I)|_{\Ker T}=-I|_{\Ker T}, \]
所以 $I|_{\Ker T}$ 是紧的，从而 $\Ker T$ 是有限维的。

类似地，另一个算子 $TS-I$ 将给出 $\coker T$ 也是有限维的，前提是我们已知 $\im T$ 是闭的。然后
\[ \|ST-I-F\|<1/2, \]
其中 $F$ 是有限秩的。于是若 $x\in\Ker F$，
\[ \|STx-x\|\le(1/2)\|x\|, \]

% [Tutor Note: p.59 Atkinson 定理证明中「连续逆 $S:(\im T)\to(\im T)^\perp$」的靶空间应为 $(\Ker T)^\perp$（$S$ 本应是双射 $(\Ker T)^\perp\to\im T$ 的逆）；500 dpi 复核确认原文确实印作 $(\im T)^\perp$，属原文笔误，正文照原样翻译。又：$H_0=(\Ker T)^\perp\oplus\Ker T$、$H_1=\im T\oplus(\im T)^\perp$ 是 Hilbert 空间的正交分解，原文写作 "disjoint union"。]

% [原文 p.60]
于是
\[ \|STx\|\ge(1/2)\|x\|, \]
从而
\[ \|Tx\|\ge\frac{1}{2\|S\|}x, \]
即 $T$ 在 $\Ker F$ 上有下界，从而 $T(\Ker F)$ 是闭的。但
\[ \im T=T(\Ker F)+T((\Ker F)^\perp) \]
是闭的。
\end{proof}

我们把 $H$ 到自身的所有 Fredholm 算子的集合记作 $\mathrm{Fred}(H)$。

\begin{defn}{5} 若 $\in\mathrm{Fred}(H)$，则定义 $T$ 的\emph{指标}，记作 $\mathrm{Index}\,T$（或简写为 $\ind T$），为差
\[ \ind T:=\dim(\Ker T)-\dim(\coker T)=[\Ker T]-[\coker T]\in\Z. \]
\end{defn}

显然，Fredholm 算子的指标是整数。

\emph{直观上可以说，指标衡量一个算子离可逆有多远，因为对可逆算子指标为零。}

指标的一个基本性质是：若 $T_0,T_1\in\mathrm{Fred}(H)$，则 $T_0T_1\in\mathrm{Fred}(H)$，且
\[ \ind(T_0T_1)=\ind T_0+\ind T_1. \]
上述性质的证明可由同调代数中的蛇引理（snake lemma，或称 serpent lemma）推出。下列序列是正合的：
\[ 0\to\Ker T_1\to\Ker T_0T_1\to\Ker T_0\to\coker T_1\to\coker T_0T_1\to\coker T_0\to 0. \]

\para{例：} 设 $H$ 有基 $e_0,e_1,e_2,\dots$，并设 $U:H\to H$ 由
\[ Ue_0=e_1 \]
\[ Ue_1=e_2 \]

% [Tutor Note: p.60 由 $\|STx\|\ge\frac12\|x\|$ 推出 $T$ 下方有界，正确的一步是 $\|STx\|\le\|S\|\cdot\|Tx\|$，故 $\|Tx\|\ge\frac{1}{2\|S\|}\|x\|$；原文末式漏了范数符号，印作 $\frac{1}{2\|S\|}x$（500 dpi 复核确认），正文照原样排。同页定义 5 的 "If $\in\mathrm{Fred}(H)$" 缺主语 $T$（应为 "If $T\in\mathrm{Fred}(H)$"），亦照原样翻译。另，原文 $Ind$ 与 $Index$ 两种记号混用，本译本统一为 $\ind$ 与 $\mathrm{Index}$。]

% [原文 p.61]
\[ \text{\emph{等等}} \]
于是 $\Ker U=\{0\}$ 且 $\dim(\coker U)=1$。因此 $\ind U=-1$。

由上例可见，指标在很大程度上取决于空间 $H_0,H_1$，而很少取决于它们之间的算子。

\begin{prop}{2} 若 $T_0,T_1\in\mathrm{Fred}(H)$ 可由 Fredholm 算子的连续道路 $T_t$，$t\in[0,1]$ 相连接，则
\[ \ind T_0=\ind T_1. \]
\end{prop}

\begin{proof}
我们只需证明映射 $t\mapsto\ind T_t$ 是局部常值的。考察 $t=t_0$ 附近。
\begin{enumerate}[label=\textbf{\alph*.},leftmargin=2.2em,itemsep=3pt]
  \item 也许 $T_{t_0}$ 是双射（从而可逆）。那么附近的 $T_t$ 也可逆，故在 $t=t_0$ 附近指标是常值的（实际上是零）。
  \item 也许 $T_{t_0}$ 是满射，设 $\dim(\Ker T_{t_0})=n$。考虑
\[ S_t:H\to H\oplus\C^n, \]
\[ x\mapsto(T_t x,P_{\Ker T_{t_0}}x). \]
这些算子是 Fredholm 的，且
\[ \ind S_t=\ind T_t-n. \]
但 $S_{t_0}$ 是双射，所以 $S_t$ 的指标在 $t_0$ 附近是常值的。
  \item 也许以上两种情形都不成立。那么我们取标准正交基 $e_i$，并令 $Q_n$ 为到 $H_n$ 上的投影，$H_n$ 是由 $e_{n+1},e_{n+2},\dots$ 生成的 $H$ 的子空间。我们断言
\[ Q_n T_{t_0}:H\to H_n \]
对充分大的 $n$ 是满射。于是
\[ \ind(Q_n T_{t_0})=\ind T_{t_0}+m \]
并应用上面的情形（b.）。
\end{enumerate}
\end{proof}

% [Tutor Note: p.61 页首的 "etc" 是 $Ue_k=e_{k+1}$（$k\ge2$）的省略；单向移位 $U$ 的核为零、余核由 $e_0$ 张成（维数 1），故 $\ind U=-1$。命题 2 情形 b 中 $S_t(x)=(T_tx,P_{\Ker T_{t_0}}x)$ 落到 $H\oplus\C^n$；情形 c 中 $Q_n$ 是到由 $e_{n+1},e_{n+2},\dots$ 生成的子空间的正交投影。$\ind(Q_nT_{t_0})=\ind T_{t_0}+m$ 中的 $m$ 原文未定义，照原文保留。]

% [原文 p.62]
命题 2 换言之是说：Fredholm 算子的集合是 $B(H)$ 中的开子集，而指标在 $\mathrm{Fred}(H)$ 上是局部常值的。

\begin{cor}{1} 若 $T$ 是 Fredholm 的、$K$ 是紧的，则 $T+K$ 是 Fredholm 的，并且
\[ \ind(T+K)=\ind T. \]
\end{cor}

\begin{proof}
考虑 $T_t=T+tK$ 并应用命题 2。
\end{proof}

% [Tutor Note: p.62 推论 1 的证明只写了 $T_t=T+tK$ 一句，补全为：$T$ 与 $T_t$ 模紧算子相同（因为 $tK$ 紧），由 Atkinson 定理（p.59 定理 1）$T_t$ 对每个 $t$ 都是 Fredholm 的，且 $t\mapsto T_t$ 连续，故由命题 2 指标常值；取 $t=0,1$ 即得 $\ind(T+K)=\ind T$。这里用到 Fredholm 算子集合在 $B(H)$ 中开（即命题 2 之后的说明）。]
